\section{Local conditional information with explicit error bounds}
\label{sec:locality}
Chronological Gaussian conditionals make information additive, but the exact
conditional at a selected time generally depends on the entire selected past.
We replace that past by a calendar window and bound the resulting error before
optimizing the schedule. The construction is Vecchia's approximation
\citep[Section 3.1]{Vecchia1988}; it is also a factorized sparse approximate
inverse \citep{Schafer2021,Huan2025,KaminetzWebber2026}. The results below
give explicit constants for the particular chronological windows used here.
They do not claim a new conditional factorization or a new filtering method.
Noisy selected subsets and vector conditional blocks are already included in
the general Vecchia framework \citep[Section 2]{KatzfussGuinness2021};
block Vecchia computation is also developed by \citet{Pan2025}.

\subsection{The exact complete-observation Markov case}
\label{subsec:exact-markov}
Suppose first that the error itself is a completely observed Gaussian Markov
chain, $\epsilon_t=X_t$, with
$X_t=A_tX_{t-1}+w_t$, $P_1=\Cov(X_1)\succ0$, and independent innovations
$\Cov(w_t)\succ0$. All blocks have the same dimension in this subsection.
Write $P_t=\Cov(X_t)$ and $\Phi_{ji}=A_j\cdots A_{i+1}$ for $i<j$.
The gap innovation covariance
$\Omega_{ji}=P_j-\Phi_{ji}P_i\Phi_{ji}^{\top}$ is positive definite.
For $S=\{t_1<\cdots<t_k\}$ the likelihood factors into its first marginal
and consecutive gap conditionals. Therefore
\begin{equation}
 I(S)=M_{t_1}+\sum_{a=1}^{k-1} W^+_{t_at_{a+1}},\quad
 M_i=F_i^\top P_i^{-1}F_i,\quad
 W^+_{ij}=(F_j-\Phi_{ji}F_i)^\top\Omega_{ji}^{-1}
                       (F_j-\Phi_{ji}F_i).
 \label{eq:exact-markov-arcs}
\end{equation}
Indeed, each gap residual is independent of the preceding selected errors;
the triangular residual map diagonalizes $R_{SS}$, giving
\eqref{eq:exact-markov-arcs} by congruence with $F_S$.

The associated ordered graph has vertices $0,1,\ldots,n,n+1$, all forward
arcs between candidate vertices, source arcs $(0,i)$ of weight $M_i$, and
zero-weight sink arcs $(i,n+1)$ and $(0,n+1)$. Each source--sink path is one
schedule. A nonnegative unit flow is a convex combination of paths: subtract
the minimum flow on any positive source--sink path and repeat; acyclicity
precludes a residual circulation. Thus its affine information image is exactly
$\operatorname{conv}\{J(S):S\subseteq[n]\}$, together with visit indicators
if these are retained. This is a direct sensitivity image of the known
principal-inverse hull of \citet{Lee2026}, building on \citet{Wei2023}.
It is not a new hull theorem.

For completeness, the mapping and the innovation-direction equivalence are
explicit. When all $A_t$ are nonsingular, let $T_1=I$ and
$T_i=A_i\cdots A_2$. Then
$R_{ij}=P_iT_i^{-\top}T_j^\top$ for $i\le j$, which is the block-factorizable
class of \citet{Lee2026}. With
$K_S=E_S^\top R_{SS}^{-1}E_S$, the affine map
$(z,K)\mapsto(z,J_0+F^\top KF)$ commutes with convex combinations.
For $C_{ij}=R_{ij}$, the reverse gap weight is
\[
 W^-_{ij}=(F_i-C_{ij}P_j^{-1}F_j)^\top
 (P_i-C_{ij}P_j^{-1}C_{ij}^\top)^{-1}
 (F_i-C_{ij}P_j^{-1}F_j).
\]
The two Schur-complement evaluations of the two-packet information give
\begin{equation}
 M_i+W^+_{ij}=W^-_{ij}+M_j.
 \label{eq:forward-reverse}
\end{equation}
Flow conservation cancels the vertex potentials $M_i$ on every fractional
flow, so forward source/gap weights and reverse gap/sink weights give the
same relaxation. Singular transitions are covered by the displayed innovation
formulas. To obtain the same conclusion from the factorized representation,
perturb every $A_t$ to nonsingular matrices converging to it, keeping positive
initial and innovation covariances fixed. The finitely many principal
inverses and arc coefficients converge. The common bounded flow polytope
therefore has the limiting affine image asserted above. A transition inverse
must not be used directly at the singular boundary.

Adding independent measurement noise usually destroys the observed Markov
property. For example, the three-packet covariance with diagonal $2$ and
off-diagonal $R_{ij}=2^{-|i-j|}$ has
$\Cov(\epsilon_1,\epsilon_3\mid\epsilon_2)=1/8\ne0$.
Complete latent observation and noisy observation therefore require different
history lengths. Appendix~\ref{app:markov-scope} also records the established
score-orthogonality extension of exact Markov paths beyond fixed-covariance
Gaussian mean information; that extension is not used to justify the noisy
finite-history bounds.

\subsection{Genuine local conditionals and a common error argument}
\label{subsec:local-conditionals}
Hereafter $\epsilon_t$ is a complete fixed packet of dimension $d_t$ in
\eqref{eq:observation-model}. For an integer $L\ge0$, set
\begin{align}
 \mathcal H_t(S)&=S\cap\{\max(1,t-L),\ldots,t-1\},\nonumber\\
 b_t&=R_{t\mathcal H_t}R_{\mathcal H_t\mathcal H_t}^{-1},&
 D_t&=R_{tt}-R_{t\mathcal H_t}R_{\mathcal H_t\mathcal H_t}^{-1}
                       R_{\mathcal H_t t},\nonumber\\
 Z_t&=\epsilon_t-\sum_{j\in\mathcal H_t}b_{tj}\epsilon_j,&
 G_{t,\mathcal H_t}&=F_t-\sum_{j\in\mathcal H_t}b_{tj}F_j.
 \label{eq:local-conditionals}
\end{align}
An empty history has no regression term. Let $A_L$ be the unit block
lower-triangular residual map on $S$ and $D_L=\diag(D_t:t\in S)$. Define
\begin{align}
 C_L(S)&=D_L^{-1/2}A_LR_{SS}A_L^\top D_L^{-1/2},&
 Q_L(S)&=A_L^\top D_L^{-1}A_L,\nonumber\\
 I_L(S)&=\sum_{t\in S}G_{t,\mathcal H_t}^\top D_t^{-1}G_{t,\mathcal H_t},&
 J_L(S)&=J_0+I_L(S).
 \label{eq:local-precision}
\end{align}
The product of the true densities $p(\epsilon_t\mid\epsilon_{\mathcal H_t})$
is a normalized Gaussian density with precision $Q_L$, as successive
integration in reverse chronological order shows. Its mean-parameter
information is $I_L$. It is a working-model information matrix, not the true
Fisher information or the covariance of an estimator under misspecification.
Every regression is recomputed on its actual local history. A full-history
Kalman filter with older coefficients deleted gives a different approximation.

\begin{lemma}[Residual-to-information transfer]
\label{lem:residual-transfer}
If $\norm{C_L(S)-I}_2\le\delta$, then
\begin{equation}
 (1-\delta)R_{SS}^{-1}\preceq Q_L(S)\preceq(1+\delta)R_{SS}^{-1}.
 \label{eq:precision-sandwich}
\end{equation}
For $J_0\succeq0$ this implies
$(1-\delta)J(S)\preceq J_L(S)\preceq(1+\delta)J(S)$.
If $\delta<1$, a sharper prior-aware upper matrix is
\begin{equation}
 J(S)\preceq J_0+\frac{I_L(S)}{1-\delta}.
 \label{eq:prior-aware-local}
\end{equation}
These statements require no sensitivity norm bound or positive prior eigenvalue.
\end{lemma}
\begin{proof}
Put $B=D_L^{-1/2}A_LR_{SS}^{1/2}$. The invertible square matrices
$BB^\top=C_L$ and $B^\top B=R_{SS}^{1/2}Q_LR_{SS}^{1/2}$ have identical
eigenvalues. Bounding these between $1-\delta$ and $1+\delta$ and applying
congruence proves \eqref{eq:precision-sandwich}. Congruence with $F_S$ gives
the data-information inequalities. Adding $J_0$ proves the remaining
statements, using $\delta J_0\succeq0$.
\end{proof}

The following lemma consolidates the overlapping-history calculation needed
for several models. Define, for $0<x<1$,
\begin{align}
 T_L(x)&=\frac{x^{L+1}}{1-x},&
 N_L(x)&=\sum_{h=1}^L\sum_{d=L+1-h}^L x^{h+2d}
 =\frac{x^{L+2}(1-x^L)(1-x^{L+1})}{(1-x)(1-x^2)}.
 \label{eq:tail-near}
\end{align}
The empty near sum is zero. Summing the inner geometric series, and then
$\sum_{h=1}^L(x^{2L+2-h}-x^{2L+2+h})/(1-x^2)$, proves the identity.

\begin{lemma}[Local regression row bound]
\label{lem:local-row}
Suppose, uniformly in selected histories,
$D_t\succeq d_* I$, $\norm{b_{tj}}_2\le B_*x^{t-j}$ for
$j\in\mathcal H_t$, and
$\norm{\Cov(Z_t,\epsilon_j)}_2\le K_*x^{t-j}$ for $j<t-L$,
where $d_*>0$, $B_*,K_*\ge0$, and $0<x<1$. Then
\begin{equation}
 \norm{C_L(S)-I}_2\le
 \frac{2K_*}{d_*}\frac{x^{L+1}}{1-x}
 \left[1+B_*\frac{x(1-x^L)}{1-x}\right].
 \label{eq:generic-row-bound}
\end{equation}
If a sharper far-pair bound
$\norm{\Cov(Z_t,Z_s)}_2\le K_*x^{t-s}$ is available for $t-s>L$,
the right side can instead be
$2K_*[T_L(x)+B_*N_L(x)]/d_*$.
\end{lemma}
\begin{proof}
For $s<t$, $h=t-s$, regression orthogonality and
\begin{equation}
 \Cov(Z_t,Z_s)=\Cov(Z_t,\epsilon_s)
 -\sum_{j\in\mathcal H_s}\Cov(Z_t,\epsilon_j)b_{sj}^\top
 \label{eq:residual-pair-identity}
\end{equation}
give the respective bounds
\[
 K_*x^h\left(1+B_*\sum_{d=1}^Lx^{2d}\right)\quad(h>L),\qquad
 K_*B_*\sum_{d=L+1-h}^Lx^{h+2d}\quad(1\le h\le L).
\]
For the near case, $s$ and all $j\ge t-L$ in the sum belong to
$\mathcal H_t$, so their covariances vanish; the older terms need not vanish.
Normalization by $D_t^{-1/2},D_s^{-1/2}$ costs at most $1/d_*$.
There is at most one packet at each distance on each side. The row sum is
at most $2K_*/d_*$ times
$T_L(x)(1+B_*\sum_{d=1}^Lx^{2d})+B_*N_L(x)$.
The terms multiplying $B_*$ sum to
$T_L(x)x(1-x^L)/(1-x)$, proving \eqref{eq:generic-row-bound}.
For any block matrix $E$, its scalar matrix of block norms majorizes
its action on the vector of component norms. If $E$ is symmetric, that
scalar matrix is symmetric and its maximum row sum bounds $\norm E_2$.
Apply this to $E=C_L-I$, whose diagonal blocks vanish. Replacing the far
majorant proves the final statement.
\end{proof}

For empty selections, independent packets, or $L\ge n-1$, the construction
is exact and we use $\delta_L=0$. With complete history, $Z_t$ is orthogonal
to every earlier selected observation and hence to every earlier residual;
with independent packets all regressions vanish. These exact cases should
be handled directly rather than through conservative infinite tails.

\subsection{Noisy scalar chains and complete observed blocks}
\label{subsec:noisy-markov}
Consider scalar errors $\epsilon_t=X_t+v_t$, where
$X_t=a_tX_{t-1}+w_t$. The centered Gaussian initial state, process noises
and measurement noises are mutually independent. Assume
\begin{equation}
 |a_t|\le\rho,\quad 0\le\rho<1,\qquad \Var(X_t)\le\overline P,
 \qquad r_t=\Var(v_t)\ge r_{\min}>0.
 \label{eq:scalar-promises}
\end{equation}
Transitions may be signed, marginal variances nonstationary, and process
variances zero. Let $\kappa=\overline P/(\overline P+r_{\min})$ and let
$d_*>0$ be any valid lower bound on all local innovation variances.
The choice $d_*=r_{\min}$ is always valid.

\begin{theorem}[Scalar noisy Markov bound]
\label{thm:scalar-locality}
Under \eqref{eq:scalar-promises}, for every selected subset,
\begin{equation}
 \norm{C_L(S)-I}_2\le\delta_L^{\mathrm{sc}}
 :=\frac{2\overline P}{d_*}
       \{T_L(\rho)+\kappa N_L(\rho)\}.
 \label{eq:scalar-refined-delta}
\end{equation}
If both latent marginal variance $P$ and measurement variance $r$ are
constant, the near coefficient $\kappa$ in this formula can be replaced by
$\kappa(1-\kappa)$, with $\overline P=P$ and $\kappa=P/(P+r)$.
At $\rho=0$ or $\overline P=0$ the error is zero.
\end{theorem}
\begin{proof}
Compute the fresh filter for $\mathcal H_t$ from the unconditional marginal
at its first observation. Its prediction variance is at most the
unconditional one. Each update has gain
$K_j=p_j^-/(p_j^-+r_j)\in[0,\kappa]$; the prediction coefficient on an
observation at $j$ is its gain multiplied by intervening transitions and
later factors $1-K_u$. Thus $|b_{tj}|\le\kappa\rho^{t-j}$.
For $j<t-L$, start the same fresh filter at time $j$ with the unconditional
state. New noises are independent of $\epsilon_j$; transitions multiply
the cross covariance by $a_u$, and updates by $1-K_u$. Consequently
\begin{equation}
 \Cov(Z_t,\epsilon_j)=\Phi_{tj}\Var(X_j)\alpha_t,
 \quad \alpha_t=\prod_{u\in\mathcal H_t}(1-K_u^{(t)}),
 \quad \Phi_{tj}=\prod_{u=j+1}^t a_u.
 \label{eq:old-scalar-identity}
\end{equation}
The gains $K_u^{(t)}$ are those of this particular local history. Their
values coincide whether one starts at $j$ without intervening updates or
at the first history observation, because unconditional marginal propagation
is exact. This proves the old-covariance bound with $K_*=\overline P$.

For $s<t-L$, put $e_s=X_s-\E(X_s\mid\epsilon_{\mathcal H_s})$ and
$p_s^-=\Var(e_s)$. Regression orthogonality gives
$\Cov(X_s,Z_s)=p_s^-$, since $Z_s=e_s+v_s$. All updates of the $t$-local
filter occur strictly after $s$; the same transport proves the exact identity
\begin{equation}
 \Cov(Z_t,Z_s)=\Phi_{ts}p_s^-\alpha_t,
 \qquad |\Cov(Z_t,Z_s)|\le\overline P\rho^{t-s}.
 \label{eq:far-scalar-identity}
\end{equation}
Lemma~\ref{lem:local-row} gives \eqref{eq:scalar-refined-delta}.
For constant $P,r$, every nonempty fresh history starts with gain exactly
$\kappa$. In a near pair $s\in\mathcal H_t$, so $\alpha_t\le1-\kappa$.
Every surviving old covariance in \eqref{eq:residual-pair-identity} receives
this factor through \eqref{eq:old-scalar-identity}. Its regression coefficient
still contributes $\kappa$, proving the stated refinement. The far bound
does not receive this factor because its local history can be empty.
\end{proof}

The simpler gain-only constant, useful for comparison with complete blocks,
is
\begin{equation}
 \delta_L^{\mathrm{gain}}=
 \frac{2\overline P}{d_*}\frac{\rho^{L+1}}{1-\rho}
 \left[1+\kappa\frac{\rho(1-\rho^L)}{1-\rho}\right].
 \label{eq:gain-delta}
\end{equation}
It follows directly from Lemma~\ref{lem:local-row}; the exact far identity
makes \eqref{eq:scalar-refined-delta} no larger. Replacing $\kappa$ by one
gives the older conservative constant
$2\overline P\rho^{L+1}(1-\rho^{L+1})/[d_*(1-\rho)^2]$.
Another improvement uses
\begin{equation}
 d_* =\min_t\Var(\epsilon_t\mid\epsilon_1,\ldots,\epsilon_{t-1})
 \ge r_{\min}.
 \label{eq:full-grid-floor}
\end{equation}
Conditioning on more observations cannot increase conditional covariance,
so one full-grid Kalman sweep gives a valid floor, while the gain constant
continues to use $r_{\min}$. Any numerical floor used in a certificate must
be bounded downward rigorously.
For stationary $P=r=1$ and $|a|=\rho$, the full-grid prediction recursion
is $x\mapsto1-\rho^2+\rho^2x/(1+x)$, started at $x=1$.
Its nonnegative fixed point is $\sqrt{1-\rho^2}$; monotonicity of the map
keeps every iterate above that point. Thus the explicit horizon-independent
floor $d_*=1+\sqrt{1-\rho^2}$ is valid in this special case.

\begin{theorem}[Complete blocks in the observation-noise metric]
\label{thm:block-locality}
Let $\epsilon_t=X_t+v_t$, $X_t=A_tX_{t-1}+w_t$, with complete blocks of
a common dimension and the same independence assumptions. Suppose
$V_t=\Cov(v_t)\succ0$ and
\begin{equation}
 \norm{V_t^{-1/2}A_tV_{t-1}^{1/2}}_2\le\rho,\quad 0\le\rho<1,\qquad
 V_t^{-1/2}P_tV_t^{-1/2}\preceq\overline P I.
 \label{eq:block-promises}
\end{equation}
Then \eqref{eq:scalar-refined-delta} bounds $\norm{C_L(S)-I}_2$ for every subset,
with $\kappa=\overline P/(1+\overline P)$ and $d_*=1$.
One may replace $1$ by any uniform lower eigenvalue bound on full-grid
innovation covariances in these whitened coordinates. The weaker gain-only
bound \eqref{eq:gain-delta} also remains valid. Neither constant has a
factor depending on block dimension.
If $\rho=0$ or $\overline P=0$, use the exact bound zero.
\end{theorem}
\begin{proof}
In coordinates $X'_t=V_t^{-1/2}X_t$, the measurement noise covariance is $I$
and the transition norm is at most $\rho$. A fresh prediction covariance
is at most $\overline P I$, so its gain
$K=P^-(P^-+I)^{-1}$ is symmetric and satisfies
$0\preceq K\preceq\kappa I$ and $0\preceq I-K\preceq I$.
Expanding the predictor by successive transitions and updates yields
$\norm{b_{tj}}_2\le\kappa\rho^{t-j}$ by submultiplicativity, with no
commutativity assumption. Propagating covariance with an excluded old
observation starts from a covariance of norm at most $\overline P$;
the same products give norm at most $\overline P\rho^{t-j}$.
For a far pair $s<t-L$, full observation gives
$\Cov(X'_s,Z'_s)=\Pi_s^{-,s}$, the prediction covariance for the
$s$-local history, of norm at most $\overline P$. All updates in the
$t$-local history occur after $s$. Therefore
$\Cov(Z'_t,Z'_s)=\mathcal T_{ts}^{(t)}\Pi_s^{-,s}$, where
$\mathcal T_{ts}^{(t)}$ is the ordered product of transitions and
intervening $I-K_u^{(t)}$ updates. Each factor has the norm bound already
proved, so this covariance has norm at most $\overline P\rho^{t-s}$.
This far-pair improvement requires no scalar commutativity.
Every innovation covariance is at least $I$. The sharp-far option of
Lemma~\ref{lem:local-row} proves the stronger bound, and its original
option gives \eqref{eq:gain-delta}. Full-grid conditioning gives the
alternative floor.
Whiten $F_t$ together with each packet. Local regression residuals then
transform by $V_t^{-1/2}$, so both true and local information are invariant.
Normalized residual covariances before and after whitening are related by
a block orthogonal congruence: each block
$(V_t^{-1/2}D_tV_t^{-1/2})^{-1/2}V_t^{-1/2}D_t^{1/2}$ is orthogonal.
The norm and precision conclusions therefore hold in the original coordinates.
\end{proof}

The contraction in \eqref{eq:block-promises} must be checked in the indicated
noise metric; contraction of $A_t$ in unrelated coordinates is insufficient.
This theorem permits noncommuting transitions and time-varying correlated
noise within packets. It does not allow selecting arbitrary channels from a
packet or replacing full observation by an arbitrary observation matrix.

\subsection{Spacing and finite-horizon row bounds}
\label{subsec:spacing-locality}
For a stationary scalar covariance
$R_{ij}=Pa^{|i-j|}+r\ind_{i=j}$, let $b=|a|<1$, $P\ge0$, $r>0$, and
require consecutive selected times to differ by at least $g\ge1$.
Set $\kappa=P/(P+r)$ and $m_L=\lfloor L/g\rfloor$.
Define $T_g^{\circ0}(P)=P$ and
\begin{equation}
 T_g(x)=P(1-b^{2g})+b^{2g}\frac{rx}{r+x},\qquad
 d_*^{\mathrm{sp}}=r+T_g^{\circ m_L}(P).
 \label{eq:spacing-floor}
\end{equation}
For distances $h\ge g$ put
\begin{equation}
 \phi(h)=\begin{cases}
 P b^h,&h>L,\\
 P\kappa(1-\kappa)b^h
 \displaystyle\sum_{d=\ell_h,\ell_h+g,\ldots\le L} b^{2d},&g\le h\le L,
 \end{cases}\quad \ell_h=\max(g,L+1-h).
 \label{eq:spacing-pair-majorant}
\end{equation}
For integers $M\ge0$ define
\begin{equation}
 B(M)=\begin{cases}0,&M<g,\\
 \max\{B(M-1),\phi(M)+B(M-g)\},&M\ge g.
 \end{cases}
 \label{eq:spacing-row-dp}
\end{equation}

\begin{proposition}[A finite-horizon bound respecting sampling gaps]
\label{prop:spacing-locality}
For every schedule satisfying the minimum gap,
\begin{equation}
 \norm{C_L(S)-I}_2\le\delta_{n,L,g}:=
 \max_{1\le t\le n}\frac{B(t-1)+B(n-t)}{d_*^{\mathrm{sp}}}.
 \label{eq:spacing-delta}
\end{equation}
The bound is zero directly if $P=0$, $b=0$, $L\ge n-1$, or $g\ge n$.
\end{proposition}
\begin{proof}
A fresh local history has at most $m_L$ observations. An observation with
prediction variance $x\le P$, followed by $d\ge g$ calendar steps, has next
prediction variance $T_d(x)$. This is increasing in $x$; it is nondecreasing
in $d$, because $rx/(r+x)\le P$. Iteration of the increasing map $T_g$
from $P$ decreases and remains in $[0,P]$. Starting from the unconditional
first prediction, replacing every gap by $g$ and adding observations up to
$m_L$ can only decrease the target variance. This proves
$D_t\ge d_*^{\mathrm{sp}}$.
The far bound is \eqref{eq:far-scalar-identity}. For a near pair, surviving
history distances $d=s-j$ lie in $[\ell_h,L]$ and are $g$-separated. If
their ordered values are $d_1<\cdots<d_q$, then
$d_i\ge\ell_h+(i-1)g$. Since $b^{2d}$ decreases, summing these largest
possible weights and applying the stationary near factor gives
\eqref{eq:spacing-pair-majorant}.
The largest sum of these nonnegative pair weights on $g$-separated distances
in $\{g,\ldots,M\}$ is exactly $B(M)$: an optimal set either omits $M$ or
includes it and has all earlier distances at most $M-g$. Left and right
distances from an anchor satisfy these restrictions separately. Bounding
each side and dividing by the innovation floor proves the row bound,
and the block row argument of Lemma~\ref{lem:local-row} proves the norm
bound. We do not assume that one schedule attains all pair majorants.
\end{proof}

Using the unrefined near coefficient $\kappa$ instead of
$\kappa(1-\kappa)$, or the far weight
$Pb^h[1+\kappa\sum_{j=1}^{m_L}b^{2gj}]$, also gives valid but weaker
bounds. These distinguish earlier certificates from those using
\eqref{eq:spacing-pair-majorant}. The stationary factor cannot be inferred
from only upper latent-variance and lower noise bounds; an exact three-time
counterexample appears in Appendix~\ref{app:locality-boundaries}.

\subsection{General covariance decay, including unequal packet sizes}
\label{subsec:general-decay}
The next result avoids a latent Markov assumption. Let $d_t$ be arbitrary
positive packet dimensions and suppose
\begin{equation}
 R\succeq mI,\qquad \norm{R_{ts}}_2\le C\rho^{|t-s|}\quad(t\ne s),
 \qquad m,C>0,\quad 0<\rho<1.
 \label{eq:decay-promises}
\end{equation}
No upper bound on diagonal blocks is imposed. Define
\begin{align}
 \vartheta&=\frac{4C\rho+m\rho(1-\rho)}{4C\rho+m(1-\rho)},&
 B_*&=\frac{2C}{m}\frac{\rho}{\vartheta-\rho},\nonumber\\
 K_*&=C\left(1+B_*\frac{\rho}{\vartheta-\rho}\right),&
 K_{\mathrm{dec}}&=\frac{2K_*}{m(1-\vartheta)}
            \left(1+B_*\frac{\vartheta}{1-\vartheta}\right).
 \label{eq:decay-constants}
\end{align}

\begin{theorem}[Explicit relative precision from covariance decay]
\label{thm:general-decay}
Under \eqref{eq:decay-promises}, every selected subset satisfies
\begin{equation}
 \norm{C_L(S)-I}_2\le\delta_L^{\mathrm{dec}}:=
 \frac{2K_*}{m}\frac{\vartheta^{L+1}}{1-\vartheta}
 \left[1+B_*\frac{\vartheta(1-\vartheta^L)}{1-\vartheta}\right]
 \le K_{\mathrm{dec}}\vartheta^{L+1}.
 \label{eq:decay-delta}
\end{equation}
The error bound and its envelope constant $K_{\mathrm{dec}}$ depend only
on $C/m$ and $\rho$, not on the horizon or any packet dimension.
For $C=0$ or $\rho=0$ the approximation is exact.
\end{theorem}
\begin{proof}
The choice in \eqref{eq:decay-constants} lies in $(\rho,1)$ and satisfies
\begin{equation}
 2C\left(\frac{\rho}{\vartheta-\rho}-\frac{\rho}{1-\rho}\right)=m/2.
 \label{eq:weighted-smallness}
\end{equation}
Consider any principal block matrix $R_{HH}$ and scalar block weights
$W_i=w_iI$ such that both ratios obey
$|w_i/w_j-1|\le\vartheta^{-|i-j|}-1$.
The diagonal blocks of
$E=W^{-1}R_{HH}W-R_{HH}$ vanish. Every row and column sum of its block
norm matrix is at most
\[
 2C\sum_{h\ge1}\rho^h(\vartheta^{-h}-1)=m/2.
\]
The row-and-column Schur estimate gives $\norm E_2\le m/2$:
if a nonnegative scalar matrix has maximum row sum $u$ and column sum $v$,
Cauchy--Schwarz yields $\norm M_2^2\le uv$, and application to component
norms gives the block version. Since $R_{HH}\succeq mI$,
$\norm{R_{HH}^{-1}E}_2\le1/2$. Factoring
$R_{HH}+E=R_{HH}(I+R_{HH}^{-1}E)$ and summing its Neumann series proves
\begin{equation}
 \norm{(W^{-1}R_{HH}W)^{-1}}_2\le2/m.
 \label{eq:conjugated-inverse}
\end{equation}
For example, $w_j=\vartheta^{-|j-i|}$ gives the usual inverse-block decay
$\norm{(R_{HH}^{-1})_{ij}}_2\le(2/m)\vartheta^{|i-j|}$.
This is the classical weighted-conjugation localization argument; here it
is used on each actual local regression.

On $\mathcal H_t$, take $w_j=\vartheta^{-(t-j)}$. These weights satisfy
the ratio condition. The horizontal row operator obeys
$\norm{R_{t\mathcal H_t}W}_2
 \le C\sum_{h\ge1}(\rho/\vartheta)^h
 =C\rho/(\vartheta-\rho)$, because its norm is at most the sum of
the block norms. The identity
\[
 b_tW=(R_{t\mathcal H_t}W)
                (W^{-1}R_{\mathcal H_t\mathcal H_t}W)^{-1}
\]
and \eqref{eq:conjugated-inverse} give
$\norm{b_{tj}}_2\le B_*\vartheta^{t-j}$.
For $j<t-L$, all $u\in\mathcal H_t$ are later than $j$, whence
\begin{align*}
 \norm{\Cov(Z_t,\epsilon_j)}_2
 &\le C\rho^{t-j}+B_*C\sum_{u\in\mathcal H_t}
                         \vartheta^{t-u}\rho^{u-j}\\
 &\le C\vartheta^{t-j}
             \left[1+B_*\sum_{r\ge1}(\rho/\vartheta)^r\right]
 =K_*\vartheta^{t-j}.
\end{align*}
Finally, the Schur-complement variational formula gives $D_t\succeq mI$:
for a vector $v$ on the target packet, $v^\top D_tv$ is the minimum
quadratic form of $R_{\{t\}\cup\mathcal H_t}$ over vectors with target
block $v$, and is therefore at least $m\norm v^2$.
Lemma~\ref{lem:local-row} proves \eqref{eq:decay-delta}; bounding
$1-\vartheta^L$ by one gives its last inequality.
\end{proof}

An alternative $\vartheta\in(\rho,1)$ may be used if the left side of
\eqref{eq:weighted-smallness} is at most $m/2$; the definitions of $B_*$
and $K_*$ then remain valid. The constants are conservative: at
$m=C=1$, $\rho=2/5$, they give $\vartheta=46/55$, $B_*=11/6$ and
$K_*=193/72$. The first windows making the finite bound at most $0.05$
and $0.01$ are respectively $49$ and $58$. These imply enormous mask
counts on a sufficiently long horizon. The theorem broadens covariance
scope but is not a claim of a practical improvement over the Markov bounds.

\begin{corollary}[Supplied block metrics]
\label{cor:decay-metric}
Let $V=\diag(V_t)$, with each $V_t\succ0$. The same theorem holds if
\eqref{eq:decay-promises} is replaced by
\begin{equation}
 R\succeq mV,\qquad
 \begin{pmatrix}
 C\rho^{|t-s|}V_t&R_{ts}\\R_{st}&C\rho^{|t-s|}V_s
 \end{pmatrix}\succeq0\quad(t\ne s).
 \label{eq:metric-decay-promises}
\end{equation}
The metrics may have arbitrarily large Euclidean condition numbers.
\end{corollary}
\begin{proof}
Congruence by $\diag(V_t^{-1/2},V_s^{-1/2})$ identifies the second
condition with
$\norm{V_t^{-1/2}R_{ts}V_s^{-1/2}}_2\le C\rho^{|t-s|}$:
the symmetric block matrix with diagonal $cI$ and off-diagonal $B,B^\top$
is positive semidefinite exactly when $\norm B_2\le c$, as follows by
singular-value decomposition. Apply Theorem~\ref{thm:general-decay}
to $V^{-1/2}RV^{-1/2}$. The block-orthogonal residual transformation in
Theorem~\ref{thm:block-locality} preserves the norm bound and congruence
returns the precision bound to the original coordinates.
\end{proof}
For rational input and rational $m,C,\rho,V_t$, the promises are rational
positive-semidefinite tests. All local inverses and information can be
computed in original coordinates; no irrational whitening is required.
Here ``metric'' means a supplied positive-definite matrix on each packet;
the calendar distance remains $|t-s|$.

A direct application is $\epsilon_t=H_tX_t+v_t$ with
$\norm{A_t}_2\le\rho$, $P_t\preceq\overline P I$,
$\norm{H_t}_2\le\overline H$, and $V_t\succeq mI$.
The independent-noise component gives $R\succeq mI$ and
$R_{ts}=H_tA_t\cdots A_{s+1}P_sH_s^\top$ for $t>s$ gives
$C=\overline H^2\overline P$. This covers fixed partial packets without
any process-noise floor, at the cost of the slower generic rate.

\subsection{Partial observations using covariance contraction}
\label{subsec:partial-locality}
For a more direct partial-observation bound, take
$X_t=A_tX_{t-1}+w_t$ and $\epsilon_t=H_tX_t+v_t$ with independent
Gaussian noises, $P_t=\Cov(X_t)$, process covariance $\Xi_t=\Cov(w_t)$,
and $V_t=\Cov(v_t)\succ0$. Latent and packet dimensions may be part of
the input. The matrices $H_t$ are fixed; selecting a packet acquires all
its responses. Assume, for constants $0\le\gamma<1$ and $s\ge0$,
\begin{equation}
 \Xi_t\succeq(1-\gamma^2)P_t\quad(t\ge2),\qquad
 H_tP_tH_t^\top\preceq sV_t\quad(t\ge1).
 \label{eq:intrinsic-partial-promises}
\end{equation}
The initial covariance may be singular. These promises are invariant under
invertible changes of state and observation coordinates and do not assume
that a partial-observation update is a Euclidean contraction.

\begin{theorem}[Intrinsic partial-packet bound]
\label{thm:partial-locality}
Put $\kappa=s/(1+s)$. Under \eqref{eq:intrinsic-partial-promises},
every selected subset satisfies
\begin{equation}
 \norm{C_L(S)-I}_2\le
 \delta_L^{\mathrm{part}}:=
 2\kappa\{T_L(\gamma)+\sqrt{\kappa s}\,N_L(\gamma)\}.
 \label{eq:partial-normalized-delta}
\end{equation}
Use zero for $s=0$, $\gamma=0$, or complete history. No dimension
appears in the constant.
\end{theorem}
\begin{proof}
We give the covariance transport details, including singular covariances.
For any fresh filter and any chosen update pattern, let $\Pi_t^-$ and
$\Pi_t^+$ denote its prediction and posterior state covariances.
Conditional covariance monotonicity gives
$\Pi_t^+\preceq\Pi_t^-\preceq P_t$.
At an update,
\[
 D=H\Pi^-H^\top+V,\qquad K=\Pi^-H^\top D^{-1},\qquad E=I-KH,
 \qquad \Pi^+=E\Pi^-E^\top+KVK^\top.
\]
At a skipped time use $E=I$ and $\Pi^+=\Pi^-$. For every transition,
\begin{equation}
 A_t\Pi_{t-1}^+A_t^\top=\Pi_t^- -\Xi_t
 \preceq\gamma^2\Pi_t^-.
 \label{eq:covariance-contraction}
\end{equation}
Indeed $\Xi_t\succeq(1-\gamma^2)P_t
\succeq(1-\gamma^2)\Pi_t^-$. Multiplying
\eqref{eq:covariance-contraction} by later update and transition matrices
and using $E\Pi^-E^\top\preceq\Pi^+$ proves
\begin{equation}
 \mathcal T_{tj}\Pi_j^+\mathcal T_{tj}^\top
       \preceq\gamma^{2(t-j)}\Pi_t^-.
 \label{eq:transport}
\end{equation}
Here $\mathcal T_{tj}$ transports a state error just after $j$ to the
prediction at $t$, including intervening updates and no update at $t$.
If there is no update at $j$, its starting covariance is the prediction
there; in particular the formula holds when a fresh filter starts with
unconditional covariance $P_j$.

Two endpoint estimates retain the signal-to-noise ratio. For any
$0\preceq\Pi\preceq P_t$, put $D=H_t\Pi H_t^\top+V_t$ and
$M=\Pi^{1/2}H_t^\top V_t^{-1}H_t\Pi^{1/2}$.
The nonzero eigenvalues of $M$ are those of
$V_t^{-1/2}H_t\Pi H_t^\top V_t^{-1/2}$, so $0\preceq M\preceq sI$.
The singular-value decomposition and the scalar inequality
$u/(1+u)\le\kappa$ for $0\le u\le s$ give
\begin{align}
 \norm{D^{-1/2}H_t\Pi^{1/2}}_2&\le\sqrt\kappa,
 \label{eq:partial-endpoint}\\
 K V_tK^\top&\preceq\kappa\Pi^+.
 \label{eq:partial-gain}
\end{align}
For the second inequality use the identities, valid also for singular $\Pi$,
$\Pi^+=\Pi^{1/2}(I+M)^{-1}\Pi^{1/2}$ and
$KV_tK^\top=\Pi^{1/2}M(I+M)^{-2}\Pi^{1/2}$.
For use below, if $UU^\top\preceq c\Pi$ then
$U=\sqrt c\Pi^{1/2}B$ for a contraction $B$ (and $U=0$ if $c=0$).
To prove this without assuming $\Pi$ invertible, diagonalize $\Pi$:
the inequality makes the rows of $U$ on its nullspace zero; divide the
remaining rows by positive square roots and use $BB^\top\preceq I$.
Together with \eqref{eq:transport} and \eqref{eq:partial-endpoint}, this
factorization converts covariance order into the following norm estimates.

For a far pair $s<t-L$, let $\Pi_s^{-,s}$ be the prediction covariance
from $\mathcal H_s$. Orthogonality yields
$U=\Cov(X_s,Z_s)=\Pi_s^{-,s}H_s^\top$ and
\[
 U D_s^{-1}U^\top\preceq\kappa\Pi_s^{-,s}\preceq\kappa P_s.
\]
The $t$-local filter starts at $s$ with $P_s$ and updates only later.
Its cross covariance with $Z_s$ is $H_t\mathcal T_{ts}U$.
Transport, factorization, and the target endpoint bound therefore give
\begin{equation}
 \norm{D_t^{-1/2}\Cov(Z_t,Z_s)D_s^{-1/2}}_2
 \le\kappa\gamma^{t-s}.
 \label{eq:partial-far}
\end{equation}
For $j<t-L$, $\Cov(X_j,\epsilon_j)=P_jH_j^\top$, and the signal promise
implies
$P_jH_j^\top V_j^{-1}H_jP_j\preceq sP_j$.
The same transport now gives
\begin{equation}
 \norm{D_t^{-1/2}\Cov(Z_t,\epsilon_j)V_j^{-1/2}}_2
 \le\sqrt{\kappa s}\,\gamma^{t-j}.
 \label{eq:partial-old}
\end{equation}
Finally, for $j\in\mathcal H_s$, its local regression coefficient is
$b_{sj}=H_s\mathcal T_{sj}K_j$ with gains from the $s$-local filter.
Using \eqref{eq:partial-gain} before transport and
\eqref{eq:partial-endpoint} after it yields
\begin{equation}
 \norm{D_s^{-1/2}b_{sj}V_j^{1/2}}_2\le\kappa\gamma^{s-j}.
 \label{eq:partial-coefficient}
\end{equation}
For a near pair, insert $V_j^{-1/2}V_j^{1/2}$ in each surviving term of
\eqref{eq:residual-pair-identity} and use
\eqref{eq:partial-old}--\eqref{eq:partial-coefficient}. Its normalized
block norm is at most
$\kappa\sqrt{\kappa s}\sum_{d=L+1-h}^L\gamma^{h+2d}$.
Summing these near weights and the far weights
\eqref{eq:partial-far} on both sides proves
\eqref{eq:partial-normalized-delta}.
If $s=0$ the observed state signal vanishes. If $\gamma=0$,
$A_tP_{t-1}A_t^\top=P_t-\Xi_t=0$, so every cross-time observation
covariance vanishes. These cases are exact even with singular state
covariances.
\end{proof}

\begin{corollary}[Euclidean process-noise floor]
\label{cor:partial-euclidean}
Suppose $P_t\preceq\overline P I$, $\Xi_t\succeq qI$ for $t\ge2$,
$V_t\succeq rI$, and $\norm{H_t}_2\le\overline H$, with
$0<q\le\overline P$ and $r>0$. Then
Theorem~\ref{thm:partial-locality} applies with
$\gamma=\sqrt{1-q/\overline P}$ and $s=\overline H^2\overline P/r$.
The earlier, weaker unnormalized bound is also valid:
\begin{equation}
 \delta_L^{\mathrm{Euc}}=
 \frac{2\overline H^2\overline P}{r}
 \left\{T_L(\gamma)+\overline H\sqrt{\overline P/r}\,N_L(\gamma)\right\}.
 \label{eq:partial-unnormalized}
\end{equation}
\end{corollary}
\begin{proof}
$qI\succeq(q/\overline P)P_t$ and
$H_tP_tH_t^\top\preceq\overline H^2\overline P I\preceq sV_t$
give the intrinsic promises. For the older constants, use
$KVK^\top\preceq\Pi^+$ in place of
\eqref{eq:partial-gain}; transport gives
$\norm{b_{tj}}_2\le\overline H\sqrt{\overline P/r}\gamma^{t-j}$.
For an old observation, $P_jH_j^\top H_jP_j
\preceq\overline H^2\overline P P_j$, so transport gives
$\norm{\Cov(Z_t,\epsilon_j)}_2
\le\overline H^2\overline P\gamma^{t-j}$.
For a far residual use
$UU^\top\preceq\overline H^2\overline P\Pi_s^{-,s}
\preceq\overline H^2\overline P P_s$ in the same way.
Lemma~\ref{lem:local-row} with $d_*=r$ proves
\eqref{eq:partial-unnormalized}. The intrinsic bound is no larger because
$\kappa\le s$ and $\kappa\sqrt{\kappa s}\le s\sqrt s$.
\end{proof}

Covariance contraction and finite-memory filtering have substantial prior
history. In particular, \citet[Theorems 1--2]{Kozdoba2019} use process-noise
dissipation in a covariance norm to approximate forecasts after burn-in for
observable time-invariant systems. Here the sufficient covariance inequalities
are imposed on the supplied finite, possibly time-varying horizon, and
transport is proved separately for every fresh local selection history.
The quantity bounded is the relative precision of the resulting local
conditionals. The process-noise floor is a substantive alternative to the
complete-block theorem, which allows zero process noise.

\subsection{A finite graph for local information and its scope}
\label{subsec:calendar-graph}
Before deciding at time $t$, a state records the last $L$ selection bits.
A skip arc shifts in zero and has information weight zero. A choose arc
shifts in one and has positive-semidefinite weight
\begin{equation}
 W_{t,H}=G_{t,H}^\top D_{t,H}^{-1}G_{t,H},
 \qquad H=\{j\in[t-L,t-1]:\text{the state selects }j\}.
 \label{eq:calendar-arc}
\end{equation}
Start from the all-zero history and connect all terminal states to a sink.
There are $O(n2^L)$ states and arcs. Each schedule yields one path and
its arc sum is exactly $I_L(S)$, because the state determines the entire
local conditional in \eqref{eq:local-conditionals}. Conversely, every path
specifies that schedule. Path decomposition therefore gives the exact hull
of the surrogate information and calendar visit vectors. If visit marginals
are binary, every path in a flow decomposition has those same decisions,
and the deterministic state transitions make the path unique.
An exact cardinality adds a count layer, giving
$O(n(k+1)2^L)$ states and arcs. Mandatory or forbidden times remove arcs.
General linking constraints preserve integer exactness, but intersection
with such constraints need not preserve the unconstrained continuous hull.

For minimum gap $g$, a valid $L$-bit mask contains no two ones within
distance $g-1$. The number of such masks is
\begin{equation}
 M(L,g)=1+\sum_{q=1}^{\lceil L/g\rceil}
             \binom{L-(g-1)(q-1)}{q}.
 \label{eq:separated-mask-count}
\end{equation}
To count masks with $q$ ones, subtract $(g-1)(i-1)$ from the $i$th ordered
one-position; this bijects to $q$ unconstrained positions in a line of
length $L-(g-1)(q-1)$. If $L<g-1$, the information mask forgets a last
observation before the spacing restriction expires. A correct graph also
stores a cooldown counter, or uses $\max(L,g-1)$ bits for feasibility
while still using only $L$ bits for information. This covers $L=0$ and
$g>L$. Feasible exact counts satisfy $k\le\lceil n/g\rceil$.
Any support objective $\sum\ip{N}{W_{t,H}}$, with fixed symmetric $N$,
is optimized exactly by the ordinary longest-path recurrence on this acyclic
graph; scalar weights need not be nonnegative for that recurrence. Complexity
and approximation guarantees for particular objectives are developed separately.

The construction uses a \emph{calendar} window. Keeping the last $m$
\emph{selected} observations instead is a valid conditional approximation,
and the scalar contraction bound applies after reindexing the selected
sequence (each selected gap contracts by at most $\rho$). A direct graph
then remembers an ordered tuple of up to $m$ calendar indices and may require
$O(n^{m+1})$ states and arcs. These state representations are different.
Likewise, refining a physical-time grid can make the per-index contraction
arbitrarily close to one. The window and graph bounds are conditional on
the stated decay constants; they do not establish uniform scalability under
grid refinement. Simultaneous scalar responses should be treated as one
packet, not as successive strictly contracting calendar steps.

\subsection{Relationship to existing locality guarantees}
\label{subsec:locality-prior}
The preceding arguments supply explicit sufficient error bounds for the
specified local histories. Their ingredients have extensive antecedents.
Inverse-factor localization is studied by \citet{BenziTuma2000},
\citet{Krishtal2015}, and \citet{Rubensson2021}; finite-predictor tail
inequalities include \citet{Meyer2015} and \citet{Inoue2018}.
The former works concern banded inverse factors, inverse-closed decay
algebras, or localized iterative factorizations; the latter use stationary
spectral or multivariate long-memory assumptions. We prove the needed
local-regression statement directly rather than presume that a theorem for
an arbitrary localized inverse factor identifies the factor in
\eqref{eq:local-precision}. In particular, the fixed exponential rate does
not follow merely from the admissible-weight algebra in
\citet[Corollary 5.1]{Krishtal2015}; its admissible weights satisfy the
Gelfand--Raikov--Shilov condition, which a fixed exponential weight fails.
Nor is an unproved direct-sum/factorization-inheritance argument used here.

\citet[Theorem 2.1]{Schafer2021} characterize fixed-pattern KL-optimal
inverse-Cholesky factors. Their Theorem 3.4 gives rigorous error and complexity
bounds for Green's covariances with geometric ordering; the neighborhood
radius has a $\log(n/\epsilon)$ dependence. Their added-noise construction
in Section 4.1 uses a latent precision approximation and incomplete Cholesky
refinement. \citet{Huan2025} study conditional selection for sparse inverse
factors, including its experimental-design interpretation.
\citet{KaminetzWebber2026} identify a partial-Cholesky-plus-Vecchia model
with an augmented Vecchia pattern and prove fixed-pattern Kaporin optimality.
These are direct antecedents for the approximation and its analysis.
The comparisons concern model hypotheses, factor and ordering, constants,
and horizon dependence; they are not a claim that KL cannot control Fisher
information. Appendix~\ref{app:locality-metrics} proves that a global Gaussian
KL bound does imply a relative precision and mean-information sandwich.

Two reductions further limit a novelty claim for
Theorem~\ref{thm:general-decay}. First, pad a selected covariance at its
original calendar positions with independent $mI$ dummy blocks at deleted
times. The padded covariance preserves \eqref{eq:decay-promises}; its
selected local conditionals are unchanged, and its normalized residual
covariance is, after permutation, $C_L(S)\oplus I$. Thus a suitable
\emph{class-uniform full-calendar} theorem already implies subset uniformity.
Second, unbounded original diagonal blocks do not preclude reduction to a
bounded-spectrum class. Put $D=\diag(R_{tt})$, $E=R-D$, and
$b_0=2C\rho/(1-\rho)$. The block row bound gives $\norm E_2\le b_0$,
while $D\succeq mI$. Hence
\[
 D\preceq R+b_0I\preceq(1+b_0/m)R,
 \qquad R\preceq D+b_0I\preceq(1+b_0/m)D,
\]
and therefore
\begin{equation}
 \frac{m}{m+b_0}I\preceq D^{-1/2}RD^{-1/2}
            \preceq(1+b_0/m)I.
 \label{eq:diagonal-priority-reduction}
\end{equation}
Its off-diagonal blocks decay with amplitude $C/m$ and rate $\rho$.
Neither subset uniformity nor absence of an upper bound on the original
diagonal alone is consequently a new localization principle. The precise
contribution supplied here is the fully explicit local-regression bounds
\eqref{eq:scalar-refined-delta}, \eqref{eq:spacing-delta},
\eqref{eq:decay-delta}, and \eqref{eq:partial-normalized-delta}, and their
use with a fixed-calendar information graph to certify the original
selected-covariance design objective. We make no first-localization claim.
